% Options for packages loaded elsewhere
\PassOptionsToPackage{unicode}{hyperref}
\PassOptionsToPackage{hyphens}{url}
%
\documentclass[
]{article}
\usepackage{amsmath,amssymb}
\usepackage{iftex}
\ifPDFTeX
  \usepackage[T1]{fontenc}
  \usepackage[utf8]{inputenc}
  \usepackage{textcomp} % provide euro and other symbols
\else % if luatex or xetex
  \usepackage{unicode-math} % this also loads fontspec
  \defaultfontfeatures{Scale=MatchLowercase}
  \defaultfontfeatures[\rmfamily]{Ligatures=TeX,Scale=1}
\fi
\usepackage{lmodern}
\ifPDFTeX\else
  % xetex/luatex font selection
\fi
% Use upquote if available, for straight quotes in verbatim environments
\IfFileExists{upquote.sty}{\usepackage{upquote}}{}
\IfFileExists{microtype.sty}{% use microtype if available
  \usepackage[]{microtype}
  \UseMicrotypeSet[protrusion]{basicmath} % disable protrusion for tt fonts
}{}
\makeatletter
\@ifundefined{KOMAClassName}{% if non-KOMA class
  \IfFileExists{parskip.sty}{%
    \usepackage{parskip}
  }{% else
    \setlength{\parindent}{0pt}
    \setlength{\parskip}{6pt plus 2pt minus 1pt}}
}{% if KOMA class
  \KOMAoptions{parskip=half}}
\makeatother
\usepackage{xcolor}
\usepackage[margin=1in]{geometry}
\usepackage{color}
\usepackage{fancyvrb}
\newcommand{\VerbBar}{|}
\newcommand{\VERB}{\Verb[commandchars=\\\{\}]}
\DefineVerbatimEnvironment{Highlighting}{Verbatim}{commandchars=\\\{\}}
% Add ',fontsize=\small' for more characters per line
\usepackage{framed}
\definecolor{shadecolor}{RGB}{248,248,248}
\newenvironment{Shaded}{\begin{snugshade}}{\end{snugshade}}
\newcommand{\AlertTok}[1]{\textcolor[rgb]{0.94,0.16,0.16}{#1}}
\newcommand{\AnnotationTok}[1]{\textcolor[rgb]{0.56,0.35,0.01}{\textbf{\textit{#1}}}}
\newcommand{\AttributeTok}[1]{\textcolor[rgb]{0.13,0.29,0.53}{#1}}
\newcommand{\BaseNTok}[1]{\textcolor[rgb]{0.00,0.00,0.81}{#1}}
\newcommand{\BuiltInTok}[1]{#1}
\newcommand{\CharTok}[1]{\textcolor[rgb]{0.31,0.60,0.02}{#1}}
\newcommand{\CommentTok}[1]{\textcolor[rgb]{0.56,0.35,0.01}{\textit{#1}}}
\newcommand{\CommentVarTok}[1]{\textcolor[rgb]{0.56,0.35,0.01}{\textbf{\textit{#1}}}}
\newcommand{\ConstantTok}[1]{\textcolor[rgb]{0.56,0.35,0.01}{#1}}
\newcommand{\ControlFlowTok}[1]{\textcolor[rgb]{0.13,0.29,0.53}{\textbf{#1}}}
\newcommand{\DataTypeTok}[1]{\textcolor[rgb]{0.13,0.29,0.53}{#1}}
\newcommand{\DecValTok}[1]{\textcolor[rgb]{0.00,0.00,0.81}{#1}}
\newcommand{\DocumentationTok}[1]{\textcolor[rgb]{0.56,0.35,0.01}{\textbf{\textit{#1}}}}
\newcommand{\ErrorTok}[1]{\textcolor[rgb]{0.64,0.00,0.00}{\textbf{#1}}}
\newcommand{\ExtensionTok}[1]{#1}
\newcommand{\FloatTok}[1]{\textcolor[rgb]{0.00,0.00,0.81}{#1}}
\newcommand{\FunctionTok}[1]{\textcolor[rgb]{0.13,0.29,0.53}{\textbf{#1}}}
\newcommand{\ImportTok}[1]{#1}
\newcommand{\InformationTok}[1]{\textcolor[rgb]{0.56,0.35,0.01}{\textbf{\textit{#1}}}}
\newcommand{\KeywordTok}[1]{\textcolor[rgb]{0.13,0.29,0.53}{\textbf{#1}}}
\newcommand{\NormalTok}[1]{#1}
\newcommand{\OperatorTok}[1]{\textcolor[rgb]{0.81,0.36,0.00}{\textbf{#1}}}
\newcommand{\OtherTok}[1]{\textcolor[rgb]{0.56,0.35,0.01}{#1}}
\newcommand{\PreprocessorTok}[1]{\textcolor[rgb]{0.56,0.35,0.01}{\textit{#1}}}
\newcommand{\RegionMarkerTok}[1]{#1}
\newcommand{\SpecialCharTok}[1]{\textcolor[rgb]{0.81,0.36,0.00}{\textbf{#1}}}
\newcommand{\SpecialStringTok}[1]{\textcolor[rgb]{0.31,0.60,0.02}{#1}}
\newcommand{\StringTok}[1]{\textcolor[rgb]{0.31,0.60,0.02}{#1}}
\newcommand{\VariableTok}[1]{\textcolor[rgb]{0.00,0.00,0.00}{#1}}
\newcommand{\VerbatimStringTok}[1]{\textcolor[rgb]{0.31,0.60,0.02}{#1}}
\newcommand{\WarningTok}[1]{\textcolor[rgb]{0.56,0.35,0.01}{\textbf{\textit{#1}}}}
\usepackage{graphicx}
\makeatletter
\def\maxwidth{\ifdim\Gin@nat@width>\linewidth\linewidth\else\Gin@nat@width\fi}
\def\maxheight{\ifdim\Gin@nat@height>\textheight\textheight\else\Gin@nat@height\fi}
\makeatother
% Scale images if necessary, so that they will not overflow the page
% margins by default, and it is still possible to overwrite the defaults
% using explicit options in \includegraphics[width, height, ...]{}
\setkeys{Gin}{width=\maxwidth,height=\maxheight,keepaspectratio}
% Set default figure placement to htbp
\makeatletter
\def\fps@figure{htbp}
\makeatother
\setlength{\emergencystretch}{3em} % prevent overfull lines
\providecommand{\tightlist}{%
  \setlength{\itemsep}{0pt}\setlength{\parskip}{0pt}}
\setcounter{secnumdepth}{-\maxdimen} % remove section numbering
\ifLuaTeX
  \usepackage{selnolig}  % disable illegal ligatures
\fi
\usepackage{bookmark}
\IfFileExists{xurl.sty}{\usepackage{xurl}}{} % add URL line breaks if available
\urlstyle{same}
\hypersetup{
  pdftitle={Part B: Synthetic Data},
  pdfauthor={Group 11},
  hidelinks,
  pdfcreator={LaTeX via pandoc}}

\title{Part B: Synthetic Data}
\author{Group 11}
\date{2026-09-21}

\begin{document}
\maketitle

\section{1. Introduction}\label{introduction}

In this part, we generate synthetic data for the proposed study. The
study investigates whether adding a human demonstration to a text
instruction affects the time an embodied AI robot needs to complete a
manipulation task.

The \textbf{independent variable} is the instruction condition: -
Text-only - Text + demonstration

The \textbf{dependent variable} is task completion time, measured in
seconds.

We generate two datasets:

\begin{enumerate}
\def\labelenumi{\arabic{enumi}.}
\tightlist
\item
  A dataset in which a true effect exists.
\item
  A dataset in which the null hypothesis is true.
\end{enumerate}

The effect of the human demonstration is controlled by an adjustable
parameter called \texttt{effect\_size}.

\section{2. Load packages}\label{load-packages}

\begin{Shaded}
\begin{Highlighting}[]
\CommentTok{\# Load the packages needed for data manipulation and visualisation}
\FunctionTok{library}\NormalTok{(tidyverse)}
\end{Highlighting}
\end{Shaded}

\section{3. Set seed and Define simulation
parameters}\label{set-seed-and-define-simulation-parameters}

\begin{Shaded}
\begin{Highlighting}[]
\CommentTok{\# Set a seed to for reproducibility}
\FunctionTok{set.seed}\NormalTok{(}\DecValTok{42}\NormalTok{)}

\CommentTok{\# Total number of robot trials}
\NormalTok{N }\OtherTok{\textless{}{-}} \DecValTok{100}

\CommentTok{\# Number of trials in each experimental condition}
\NormalTok{n\_per\_group }\OtherTok{\textless{}{-}}\NormalTok{ N }\SpecialCharTok{/} \DecValTok{2}

\CommentTok{\# Mean completion time for the text{-}only condition, in seconds}
\NormalTok{mean\_text }\OtherTok{\textless{}{-}} \DecValTok{30}

\CommentTok{\# Standard deviation of completion time within each condition, in seconds}
\NormalTok{sd\_completion }\OtherTok{\textless{}{-}} \DecValTok{5}

\CommentTok{\# True effect of adding a human demonstration}
\CommentTok{\# A negative value means that the demonstration is expected to}
\CommentTok{\# reduce completion time}
\CommentTok{\# For example, {-}3 means that the demonstration reduces the}
\CommentTok{\# population mean completion time by 3 seconds}
\NormalTok{effect\_size }\OtherTok{\textless{}{-}} \SpecialCharTok{{-}}\DecValTok{3}

\CommentTok{\# Mean completion time in the demonstration condition}
\NormalTok{mean\_demo }\OtherTok{\textless{}{-}}\NormalTok{ mean\_text }\SpecialCharTok{+}\NormalTok{ effect\_size}
\end{Highlighting}
\end{Shaded}

\section{4.1 Simulate the dataset with a true
effect}\label{simulate-the-dataset-with-a-true-effect}

\begin{Shaded}
\begin{Highlighting}[]
\CommentTok{\# Generate completion times for robots receiving the text{-}only instruction}
\CommentTok{\# rnorm() draws random observations from a normal distribution.}
\CommentTok{\# The distribution has a mean of 30 seconds and an SD of 5 seconds}
\NormalTok{text\_effect }\OtherTok{\textless{}{-}} \FunctionTok{rnorm}\NormalTok{(}
  \AttributeTok{n =}\NormalTok{ n\_per\_group,}
  \AttributeTok{mean =}\NormalTok{ mean\_text,}
  \AttributeTok{sd =}\NormalTok{ sd\_completion}
\NormalTok{)}

\CommentTok{\# Generate completion times for robots receiving both the text instruction}
\CommentTok{\# and a human demonstration}
\CommentTok{\# The mean is lower by the specified effect size ({-}3 seconds)}
\NormalTok{demo\_effect }\OtherTok{\textless{}{-}} \FunctionTok{rnorm}\NormalTok{(}
  \AttributeTok{n =}\NormalTok{ n\_per\_group,}
  \AttributeTok{mean =}\NormalTok{ mean\_demo,}
  \AttributeTok{sd =}\NormalTok{ sd\_completion}
\NormalTok{)}

\CommentTok{\# Combine the two groups into one data frame}
\NormalTok{data\_effect }\OtherTok{\textless{}{-}} \FunctionTok{tibble}\NormalTok{(}
  \AttributeTok{condition =} \FunctionTok{c}\NormalTok{(}
    \FunctionTok{rep}\NormalTok{(}\StringTok{"Text{-}only"}\NormalTok{, n\_per\_group),}
    \FunctionTok{rep}\NormalTok{(}\StringTok{"Text + demonstration"}\NormalTok{, n\_per\_group)}
\NormalTok{  ),}
  \AttributeTok{completion\_time =} \FunctionTok{c}\NormalTok{(}
\NormalTok{    text\_effect,}
\NormalTok{    demo\_effect}
\NormalTok{  )}
\NormalTok{)}

\CommentTok{\# Completion time cannot be negative in reality}
\CommentTok{\# We therefore replace any negative simulated values with zero}
\NormalTok{data\_effect }\OtherTok{\textless{}{-}}\NormalTok{ data\_effect }\SpecialCharTok{\%\textgreater{}\%}
  \FunctionTok{mutate}\NormalTok{(}
    \AttributeTok{completion\_time =} \FunctionTok{pmax}\NormalTok{(completion\_time, }\DecValTok{0}\NormalTok{)}
\NormalTok{  )}

\CommentTok{\# Display the first few observations}
\FunctionTok{head}\NormalTok{(data\_effect)}
\end{Highlighting}
\end{Shaded}

\begin{verbatim}
## # A tibble: 6 x 2
##   condition completion_time
##   <chr>               <dbl>
## 1 Text-only            36.9
## 2 Text-only            27.2
## 3 Text-only            31.8
## 4 Text-only            33.2
## 5 Text-only            32.0
## 6 Text-only            29.5
\end{verbatim}

\section{4.2 Inspect the simulated effect
dataset}\label{inspect-the-simulated-effect-dataset}

\begin{Shaded}
\begin{Highlighting}[]
\CommentTok{\# Calculate descriptive statistics for each experimental condition}
\NormalTok{data\_effect }\SpecialCharTok{\%\textgreater{}\%}
  \FunctionTok{group\_by}\NormalTok{(condition) }\SpecialCharTok{\%\textgreater{}\%}
  \FunctionTok{summarise}\NormalTok{(}
    \AttributeTok{n =} \FunctionTok{n}\NormalTok{(),}
    \AttributeTok{mean =} \FunctionTok{mean}\NormalTok{(completion\_time),}
    \AttributeTok{sd =} \FunctionTok{sd}\NormalTok{(completion\_time),}
    \AttributeTok{min =} \FunctionTok{min}\NormalTok{(completion\_time),}
    \AttributeTok{max =} \FunctionTok{max}\NormalTok{(completion\_time)}
\NormalTok{  )}
\end{Highlighting}
\end{Shaded}

\begin{verbatim}
## # A tibble: 2 x 6
##   condition                n  mean    sd   min   max
##   <chr>                <int> <dbl> <dbl> <dbl> <dbl>
## 1 Text + demonstration    50  27.5  4.62  12.0  34.9
## 2 Text-only               50  29.8  5.76  16.7  41.4
\end{verbatim}

\section{5.1 Simulate the dataset under null
hypothesis}\label{simulate-the-dataset-under-null-hypothesis}

\begin{Shaded}
\begin{Highlighting}[]
\CommentTok{\# For the null dataset, we set the true effect to zero, which means that }
\CommentTok{\# the population mean completion time is identical in the two conditions}
\NormalTok{null\_effect }\OtherTok{\textless{}{-}} \DecValTok{0}

\CommentTok{\# Calculate the mean completion time for the demonstration condition}
\CommentTok{\# under the null hypothesis}
\NormalTok{mean\_demo\_null }\OtherTok{\textless{}{-}}\NormalTok{ mean\_text }\SpecialCharTok{+}\NormalTok{ null\_effect}

\CommentTok{\# Generate completion times for the text{-}only condition}
\NormalTok{text\_null }\OtherTok{\textless{}{-}} \FunctionTok{rnorm}\NormalTok{(}
  \AttributeTok{n =}\NormalTok{ n\_per\_group,}
  \AttributeTok{mean =}\NormalTok{ mean\_text,}
  \AttributeTok{sd =}\NormalTok{ sd\_completion}
\NormalTok{)}

\CommentTok{\# Generate completion times for the demonstration condition}
\CommentTok{\# Since null\_effect = 0, both groups have the same population mean}
\NormalTok{demo\_null }\OtherTok{\textless{}{-}} \FunctionTok{rnorm}\NormalTok{(}
  \AttributeTok{n =}\NormalTok{ n\_per\_group,}
  \AttributeTok{mean =}\NormalTok{ mean\_demo\_null,}
  \AttributeTok{sd =}\NormalTok{ sd\_completion}
\NormalTok{)}

\CommentTok{\# Combine the two groups into one data frame}
\NormalTok{data\_null }\OtherTok{\textless{}{-}} \FunctionTok{tibble}\NormalTok{(}
  \AttributeTok{condition =} \FunctionTok{c}\NormalTok{(}
    \FunctionTok{rep}\NormalTok{(}\StringTok{"Text{-}only"}\NormalTok{, n\_per\_group),}
    \FunctionTok{rep}\NormalTok{(}\StringTok{"Text + demonstration"}\NormalTok{, n\_per\_group)}
\NormalTok{  ),}
  \AttributeTok{completion\_time =} \FunctionTok{c}\NormalTok{(}
\NormalTok{    text\_null,}
\NormalTok{    demo\_null}
\NormalTok{  )}
\NormalTok{)}

\CommentTok{\# Ensure that completion times cannot be negative}
\NormalTok{data\_null }\OtherTok{\textless{}{-}}\NormalTok{ data\_null }\SpecialCharTok{\%\textgreater{}\%}
  \FunctionTok{mutate}\NormalTok{(}
    \AttributeTok{completion\_time =} \FunctionTok{pmax}\NormalTok{(completion\_time, }\DecValTok{0}\NormalTok{)}
\NormalTok{  )}

\CommentTok{\# Display the first few observations}
\FunctionTok{head}\NormalTok{(data\_null)}
\end{Highlighting}
\end{Shaded}

\begin{verbatim}
## # A tibble: 6 x 2
##   condition completion_time
##   <chr>               <dbl>
## 1 Text-only            36.0
## 2 Text-only            35.2
## 3 Text-only            25.0
## 4 Text-only            39.2
## 5 Text-only            26.7
## 6 Text-only            30.5
\end{verbatim}

\section{5.2 Inspect the null dataset}\label{inspect-the-null-dataset}

\begin{Shaded}
\begin{Highlighting}[]
\CommentTok{\# Calculate descriptive statistics for each condition:}
\CommentTok{\# Even though the population means are identical, the sample means}
\CommentTok{\# will usually differ slightly because of random sampling variation}
\NormalTok{data\_null }\SpecialCharTok{\%\textgreater{}\%}
  \FunctionTok{group\_by}\NormalTok{(condition) }\SpecialCharTok{\%\textgreater{}\%}
  \FunctionTok{summarise}\NormalTok{(}
    \AttributeTok{n =} \FunctionTok{n}\NormalTok{(),}
    \AttributeTok{mean =} \FunctionTok{mean}\NormalTok{(completion\_time),}
    \AttributeTok{sd =} \FunctionTok{sd}\NormalTok{(completion\_time),}
    \AttributeTok{min =} \FunctionTok{min}\NormalTok{(completion\_time),}
    \AttributeTok{max =} \FunctionTok{max}\NormalTok{(completion\_time)}
\NormalTok{  )}
\end{Highlighting}
\end{Shaded}

\begin{verbatim}
## # A tibble: 2 x 6
##   condition                n  mean    sd   min   max
##   <chr>                <int> <dbl> <dbl> <dbl> <dbl>
## 1 Text + demonstration    50  29.9  4.42  21.3  39.1
## 2 Text-only               50  29.2  4.64  19.9  43.5
\end{verbatim}

\section{6. Visualize the dataset with a true
effect}\label{visualize-the-dataset-with-a-true-effect}

\begin{Shaded}
\begin{Highlighting}[]
\CommentTok{\# Plot the distribution of completion times in the two conditions:}

\CommentTok{\# Individual observations are shown as points}
\CommentTok{\# the boxplots summarise the distribution within each condition.}
\FunctionTok{ggplot}\NormalTok{(}
\NormalTok{  data\_effect,}
  \FunctionTok{aes}\NormalTok{(}
    \AttributeTok{x =}\NormalTok{ condition,}
    \AttributeTok{y =}\NormalTok{ completion\_time}
\NormalTok{  )}
\NormalTok{) }\SpecialCharTok{+}
  \FunctionTok{geom\_boxplot}\NormalTok{(}\AttributeTok{outlier.shape =} \ConstantTok{NA}\NormalTok{) }\SpecialCharTok{+}
  \FunctionTok{geom\_jitter}\NormalTok{(}\AttributeTok{width =} \FloatTok{0.15}\NormalTok{, }\AttributeTok{alpha =} \FloatTok{0.6}\NormalTok{) }\SpecialCharTok{+}
  \FunctionTok{labs}\NormalTok{(}
    \AttributeTok{title =} \StringTok{"Simulated Completion Time: True Effect"}\NormalTok{,}
    \AttributeTok{x =} \StringTok{"Instruction condition"}\NormalTok{,}
    \AttributeTok{y =} \StringTok{"Completion time (seconds)"}
\NormalTok{  ) }\SpecialCharTok{+}
  \FunctionTok{theme\_minimal}\NormalTok{()}
\end{Highlighting}
\end{Shaded}

\includegraphics{a1_files/figure-latex/unnamed-chunk-6-1.pdf}

\section{7. Visualize the null
dataset}\label{visualize-the-null-dataset}

\begin{Shaded}
\begin{Highlighting}[]
\CommentTok{\# Visualise the completion times under the null hypothesis:}
\CommentTok{\# The two groups should overlap substantially because their population}
\CommentTok{\# means are identical.}
\FunctionTok{ggplot}\NormalTok{(}
\NormalTok{  data\_null,}
  \FunctionTok{aes}\NormalTok{(}
    \AttributeTok{x =}\NormalTok{ condition,}
    \AttributeTok{y =}\NormalTok{ completion\_time}
\NormalTok{  )}
\NormalTok{) }\SpecialCharTok{+}
  \FunctionTok{geom\_boxplot}\NormalTok{(}\AttributeTok{outlier.shape =} \ConstantTok{NA}\NormalTok{) }\SpecialCharTok{+}
  \FunctionTok{geom\_jitter}\NormalTok{(}\AttributeTok{width =} \FloatTok{0.15}\NormalTok{, }\AttributeTok{alpha =} \FloatTok{0.6}\NormalTok{) }\SpecialCharTok{+}
  \FunctionTok{labs}\NormalTok{(}
    \AttributeTok{title =} \StringTok{"Simulated Completion Time: Null Effect"}\NormalTok{,}
    \AttributeTok{x =} \StringTok{"Instruction condition"}\NormalTok{,}
    \AttributeTok{y =} \StringTok{"Completion time (seconds)"}
\NormalTok{  ) }\SpecialCharTok{+}
  \FunctionTok{theme\_minimal}\NormalTok{()}
\end{Highlighting}
\end{Shaded}

\includegraphics{a1_files/figure-latex/unnamed-chunk-7-1.pdf}

\section{8. Compare the two conditions
statistically}\label{compare-the-two-conditions-statistically}

\begin{Shaded}
\begin{Highlighting}[]
\CommentTok{\# Perform the planned independent{-}samples two{-}sided t{-}test}
\CommentTok{\# on the dataset containing a true effect:}
\CommentTok{\# The test compares mean completion time between the two}
\CommentTok{\# instruction conditions}
\NormalTok{t\_test\_effect }\OtherTok{\textless{}{-}} \FunctionTok{t.test}\NormalTok{(}
\NormalTok{  completion\_time }\SpecialCharTok{\textasciitilde{}}\NormalTok{ condition,}
  \AttributeTok{data =}\NormalTok{ data\_effect,}
  \AttributeTok{alternative =} \StringTok{"two.sided"}
\NormalTok{)}

\CommentTok{\# Display the results.}
\NormalTok{t\_test\_effect}
\end{Highlighting}
\end{Shaded}

\begin{verbatim}
## 
##  Welch Two Sample t-test
## 
## data:  completion_time by condition
## t = -2.2196, df = 93.647, p-value = 0.02886
## alternative hypothesis: true difference in means between group Text + demonstration and group Text-only is not equal to 0
## 95 percent confidence interval:
##  -4.3918920 -0.2443761
## sample estimates:
## mean in group Text + demonstration            mean in group Text-only 
##                           27.50351                           29.82164
\end{verbatim}

\begin{Shaded}
\begin{Highlighting}[]
\CommentTok{\# Perform the same planned two{-}sided t{-}test on the null dataset:}
\CommentTok{\# Here the true population effect is zero, so any observed}
\CommentTok{\# difference between the sample means is due to random variation.}
\NormalTok{t\_test\_null }\OtherTok{\textless{}{-}} \FunctionTok{t.test}\NormalTok{(}
\NormalTok{  completion\_time }\SpecialCharTok{\textasciitilde{}}\NormalTok{ condition,}
  \AttributeTok{data =}\NormalTok{ data\_null,}
  \AttributeTok{alternative =} \StringTok{"two.sided"}
\NormalTok{)}

\CommentTok{\# Display the results.}
\NormalTok{t\_test\_null}
\end{Highlighting}
\end{Shaded}

\begin{verbatim}
## 
##  Welch Two Sample t-test
## 
## data:  completion_time by condition
## t = 0.70345, df = 97.783, p-value = 0.4834
## alternative hypothesis: true difference in means between group Text + demonstration and group Text-only is not equal to 0
## 95 percent confidence interval:
##  -1.161274  2.436622
## sample estimates:
## mean in group Text + demonstration            mean in group Text-only 
##                           29.88142                           29.24374
\end{verbatim}

\end{document}
